\chapter{Reference: focus and tab order}
\label{ch:refman-focus}

\begin{aside}
When an app has multiple focusable regions (form fields,
panels, a sidebar), one of them is "focused" at any moment.
Track that in the Model and route keys accordingly.
\end{aside}

\section{Working example}

\begin{lstlisting}[language=C++]
#include <maya/maya.hpp>

using namespace maya;
using namespace maya::dsl;

struct Focus {
    enum class Panel { Left, Right };
    struct Model {
        Panel which = Panel::Left;
        int left = 0;
        int right = 0;
    };
    struct Switch {};
    struct Bump {};
    struct Quit {};
    using Msg = std::variant<Switch, Bump, Quit>;

    static Model init() { return {}; }

    static auto update(Model m, Msg msg)
        -> std::pair<Model, Cmd<Msg>>
    {
        return std::visit(overload{
            [&](Switch) {
                m.which = (m.which == Panel::Left)
                    ? Panel::Right : Panel::Left;
                return std::pair{m, Cmd<Msg>{}};
            },
            [&](Bump) {
                if (m.which == Panel::Left) m.left++;
                else                        m.right++;
                return std::pair{m, Cmd<Msg>{}};
            },
            [](Quit) {
                return std::pair{Model{}, Cmd<Msg>::quit()};
            },
        }, msg);
    }

    static Element view(const Model& m) {
        auto panel = [&](std::string name, int v, bool active) {
            auto body = v(
                text(name) | Bold,
                blank_,
                text(std::to_string(v))
            ) | pad<1>;
            return active
                ? body | border_<Heavy>
                : body | border_<Round> | Dim;
        };

        return v(
            h(
                panel("left",  m.left,  m.which == Panel::Left)  | flex<1>,
                panel("right", m.right, m.which == Panel::Right) | flex<1>
            ),
            blank_,
            t<"tab switch, space bump, q quit"> | Dim
        );
    }

    static auto subscribe(const Model&) -> Sub<Msg> {
        return key_map<Msg>({
            {'q',             Quit{}},
            {SpecialKey::Tab, Switch{}},
            {' ',             Bump{}},
        });
    }
};

int main() { run<Focus>({.title = "focus"}); }
\end{lstlisting}

\section{Notes}

\begin{itemize}[leftmargin=*]
  \item Make the active panel visually obvious; heavy border
    plus undimmed text beats subtle marks.
  \item For many panels, store a generic index and a fixed
    array of panel handlers.
  \item Always include a visible hint for the focus key.
\end{itemize}

\section{Recap}

\begin{recap}
\begin{itemize}[leftmargin=*]
  \item One enum (or index) in the Model.
  \item Tab cycles; render distinguishes.
  \item Active panel handles the contextual keys.
\end{itemize}
\end{recap}

\section*{Workshop}

\begin{enumerate}[leftmargin=*]
  \item Add a third panel and adjust the cycle.
  \item Use \code{shift-tab} for backwards cycle.
  \item Make the panels show different content per state.
\end{enumerate}
